\chapter{Estado del arte}

Esta fase se apoya en varias líneas de trabajo recientes: recuperación
aumentada con conocimiento estructurado, sistemas multi-agente basados en
modelos de lenguaje, uso de herramientas, auto-revisión y desarrollo
asistido por \textit{coding agents}. Este capítulo resume las ideas que se
incorporan al sistema y cierra situando el trabajo frente a ellas.

\section[Recuperación aumentada y memoria]{Recuperación aumentada y memoria de agentes}

La recuperación aumentada con conocimiento externo (RAG), formulada por
Lewis et al. \cite{lewis2020rag}, estableció un patrón hoy estándar en
sistemas con modelos de lenguaje: en lugar de depender exclusivamente
del contexto contenido en el modelo o del system prompts, el sistema recupera fragmentos
relevantes de un corpus externo y los incorpora a la generación. Dense
Passage Retrieval \cite{karpukhin2020dpr} consolidó la recuperación
densa mediante \textit{embeddings} asimétricos, mientras que BM25
\cite{robertson2009bm25} sigue siendo la base de la recuperación léxica
con la que la densa se combina.



Aproximaciones recientes añaden estructura adicional al corpus.
HippoRAG \cite{gutierrez2024hipporag} reorganiza la memoria como un
grafo de conocimiento y aplica PageRank personalizado para integrar
recuperaciones multimodales, inspirándose en una metáfora neurobiológica
de la memoria a largo plazo. Microsoft GraphRAG \cite{edge2024graphrag}
demuestra que construir un grafo a partir del corpus, combinado con
resúmenes locales y globales, permite responder preguntas que la
recuperación plana no resuelve, especialmente cuando la respuesta exige
agregar evidencia distribuida por el corpus.

Las arquitecturas de grafo de conocimiento, formalizadas en el
\textit{survey} de Hogan et al. \cite{hogan2021knowledgegraphs}, ofrecen
un sustrato para representar relaciones explícitas entre entidades.
En el laboratorio desarrollado, esa estructura encaja bien: paradigmas,
formulaciones, autores y procedencias forman relaciones que el grafo expone
directamente y que un índice plano de documentos no captura. Esta fase
combina ese grafo con dos modos complementarios de recuperación sobre el
mismo \textit{backbone}: un índice denso que aproxima la similitud
semántica entre fragmentos --- útil cuando la pregunta se formula con
otras palabras que la memoria almacenada --- y un índice léxico
BM25 que prioriza coincidencias exactas de términos --- útil cuando lo
relevante es un nombre concreto, un identificador o una cita literal. El
detalle de qué consulta utiliza cada agente y cómo se mezclan los
resultados se desarrolla en el capítulo de diseño.

\section{Sistemas multi-agente basados en LLMs}

La forma más sencilla de utilizar un modelo de lenguaje es como un agente
único: un sistema que mantiene todo el contexto, recibe la entrada del
usuario y produce la respuesta directamente. Esta aproximación funciona
bien para tareas acotadas o sencillas, pero se degrada cuando hay que coordinar
varias subtareas con dependencias, invocar herramientas externas o
validar resultados intermedios; en esos casos, el agente único tiende a
producir respuestas convincentes pero difíciles de auditar paso a paso. A
medida que el contexto se llena con instrucciones, datos y resultados
intermedios, el modelo pierde precisión sobre lo que se le ha pedido y
crece la probabilidad de que invente detalles que no estaban en la
entrada original --- el fenómeno conocido como alucinación.

Park et al. \cite{park2023generativeagents} mostraron empíricamente que
un sistema multi-agente con memoria persistente, planificación y
reflexión produce comportamientos emergentes más ricos que un agente
único operando sobre el mismo conjunto de \textit{prompts}. Sobre esa
intuición se han construido marcos de coordinación; entre ellos,
MetaGPT \cite{hong2024metagpt} es especialmente cercano al adoptado en
este trabajo, ya que codifica procedimientos operativos estandarizados
como secuencias de roles especializados (Architect, Engineer, etc.) que
se invocan en un orden canónico.

El \textit{survey} de Guo et al. \cite{guo2024llmbased} sistematiza el
campo y propone una taxonomía con tres paradigmas de comunicación ---
cooperativo, debate, competitivo --- y cuatro patrones estructurales
--- \textit{Layered}, \textit{Decentralized}, \textit{Centralized} y
\textit{Shared Message Pool}. La arquitectura adoptada en esta fase, en
la que un Orchestrator centraliza la comunicación y un conjunto de
agentes especializados reciben tareas acotadas, corresponde a una
composición \textit{Centralized} con estructura \textit{Layered}, y se
distingue por canalizar a través del Orchestrator la mayor parte del
acceso al \textit{Knowledge Backbone}: solo el Analyst lo consulta de
forma directa como \textit{tool call}, mientras que el resto de agentes
recibe un contexto pre-recuperado.

\section{Uso de herramientas y auto-revisión}

ReAct \cite{yao2023react} introdujo el patrón de intercalar razonamiento
y acción: el modelo de lenguaje produce trazas de pensamiento y llamadas
a herramientas externas, observa los resultados y razona sobre ellos
antes de actuar de nuevo. Toolformer \cite{schick2023toolformer}
demostró que los propios modelos pueden aprender, sin supervisión
exhaustiva, cuándo invocar herramientas externas. El catálogo de
herramientas del Orchestrator es una instancia directa de este patrón,
aplicado a la coordinación del pipeline de simulación, observación y
análisis.

Madaan et al. \cite{madaan2023selfrefine} y Shinn et al.
\cite{shinn2023reflexion} propusieron complementariamente que un modelo
de lenguaje puede mejorar sus propias salidas mediante refinamiento
iterativo o auto-crítica verbal. En este trabajo, la auto-revisión
no se aplica al pipeline en tiempo de ejecución, pero sí aparece de
forma intensiva en el ciclo de desarrollo --- mediante subagentes
adversariales que revisan el código producido por el agente principal
antes de fusionarlo, como se detalla en el capítulo correspondiente.

\section{Coding agents y desarrollo asistido}

Codex \cite{chen2021codex} estableció el \textit{baseline} de
generación de código a partir de descripciones en lenguaje natural,
junto con HumanEval como su \textit{benchmark} canónico. MBPP
\cite{austin2021program} introdujo un conjunto complementario de
problemas \textit{Python} sencillos. SWE-bench
\cite{jimenez2024swebench} y SWE-agent \cite{yang2024sweagent}
elevaron el listón: el primero, evaluando la capacidad de los modelos
sobre 2.294 \textit{issues} reales de GitHub que requieren coordinar
cambios en varios archivos; el segundo, definiendo una interfaz
agente-máquina pensada para que un modelo opere sobre un sistema
\textit{software} completo. En esa misma línea, ChatDev
\cite{qian2024chatdev} traslada la metáfora multi-agente al propio
desarrollo de \textit{software}: distintos agentes asumen los papeles
clásicos del equipo (analista, diseñador, programador, probador) y
cooperan para producir el código, una idea conceptualmente cercana al
reparto de responsabilidades que adopta este trabajo durante su ciclo
de desarrollo.

La evidencia empírica sobre el impacto de estas herramientas en la
productividad de los desarrolladores ha crecido en los últimos años.
Peng et al. \cite{peng2023impact} midieron, en un experimento
controlado, aceleraciones cercanas al 55\% en tareas de implementación
con GitHub Copilot. Cui et al. \cite{cui2025effects} extendieron el
análisis a tres ensayos aleatorizados a gran escala publicados en
\textit{Management Science}, situando el incremento medio de tareas
completadas de aproximadamente el 26\% sobre poblaciones de varios miles de
desarrolladores. Butler et al. \cite{butler2024dear} complementaron
estas medidas cuantitativas con un estudio cualitativo mixto que
registra los cambios en la experiencia diaria del trabajo. Estos tres
trabajos ofrecen el sustento empírico sobre el que el capítulo de
desarrollo de esta memoria construye su argumento sobre la
deconstrucción del ciclo clásico de desarrollo.

\section{Posicionamiento de esta fase}

Frente a los marcos revisados, este trabajo combina varios elementos que en
la literatura aparecen distribuidos. Toma de los sistemas multi-agente la
división explícita de responsabilidades entre roles especializados; de los
marcos RAG/KG, la mediación del acceso al conocimiento a través de un
\textit{backbone} híbrido; del uso de herramientas, la modelización del
pipeline como secuencia de \textit{tool calls}; y de la revisión
adversarial, la verificación intensiva del código durante el ciclo de
desarrollo.

La aportación de esta fase está en combinar esos elementos para un
problema concreto: la ejecución comparativa de paradigmas de toma de
decisiones humanas. Dos rasgos son especialmente propios del sistema. El
primero es la dualidad lectura-escritura del Knowledge Backbone: el sistema
consume paradigmas y postulados de la memoria compartida y, en la misma
sesión, persiste sus observaciones de simulación para que ejecuciones
futuras puedan apoyarse en ellas. El segundo es la carga dinámica, por
\textit{duck typing}, de modelos de decisión generados por otra fase del
proyecto. Esta separación permite que los dos sistemas evolucionen de forma
independiente sin que un cambio en uno obligue a reorganizar el otro.
